%! AP Statistics — Unit 4, Lesson 4.10 — Homework
\documentclass[10pt]{article}
\usepackage{apstats-article}
\usepackage{apstats-boxes}

\begin{document}

\pageheader{Unit 4, Lesson 4.10}{Homework}
\namedateperiod

\vspace{0.1in}

\begin{practicebox}
\small
\textbf{Directions:} For each scenario, (a) check BINS and state whether the situation
is a binomial setting, and (b) if yes, compute the requested probability showing the
formula setup.

\vspace{10pt}
\textbf{1.} A basketball player makes 72\% of her free throws. She shoots 10 free
throws in a game. Let $X =$ the number she makes.
\begin{enumerate}[label=(\alph*), leftmargin=2em, itemsep=8pt]
  \item Check BINS and state the distribution (or explain which condition fails).\\
        \writeline\\[1pt]\writeline\\[1pt]\writeline
  \item Find $P(X = 8)$. Show the formula and the numerical answer.\\
        \writeline\\[1pt]\writeline
\end{enumerate}

\vspace{8pt}
\textbf{2.} A survey asks 200 randomly selected adults (from a city of 500{,}000)
whether they support a new recycling program. Historically 60\% support such
programs. Let $X =$ the number who support the program.
\begin{enumerate}[label=(\alph*), leftmargin=2em, itemsep=8pt]
  \item Check BINS. Is the $10\%$ condition relevant here? Explain.\\
        \writeline\\[1pt]\writeline
  \item Find $P(X = 120)$ using \texttt{binompdf}. (No hand calculation required ---
        just set up the calculator notation and record the result.)\\
        \writeline
\end{enumerate}

\vspace{8pt}
\textbf{3.} A class has 30 students. The teacher randomly calls on 6 students without
replacement. 10 of the 30 students did not complete the homework.
Let $X =$ the number called on who did not complete the homework.
\begin{enumerate}[label=(\alph*), leftmargin=2em, itemsep=8pt]
  \item Is this a binomial setting? Which condition(s), if any, are violated?\\
        \writeline\\[1pt]\writeline
  \item Check the $10\%$ condition: is $n \le 0.10 \times N$? Show the calculation.\\
        \writeline
\end{enumerate}

\vspace{8pt}
\textbf{4.} Evaluate $\binom{6}{2}$ by hand. Then explain in one sentence what this
number represents in the context of a binomial experiment with $n = 6$ trials where
you are interested in exactly 2 successes.
\begin{enumerate}[label=(\alph*), leftmargin=2em, itemsep=8pt]
  \item $\binom{6}{2} = $ \blank{3cm}
  \item Interpretation: \writeline\\[1pt]\writeline
\end{enumerate}

\vspace{8pt}
\textbf{5.} (AP-Style Free Response) \textit{A door-to-door salesperson visits
8 houses. Based on past experience, the probability of making a sale at any given
house is 0.15. Assume houses are independent.}
\begin{enumerate}[label=(\alph*), leftmargin=2em, itemsep=8pt]
  \item Define the random variable $X$ and state its distribution. Check BINS
        briefly.\\
        \writeline\\[1pt]\writeline
  \item Calculate $P(X = 1)$. Show formula setup.\\
        \writeline\\[1pt]\writeline
  \item Calculate $P(X = 0)$. Interpret in context.\\
        \writeline\\[1pt]\writeline
\end{enumerate}

\vspace{8pt}
\textbf{6. Extension:} Build the complete probability distribution table for
$X \sim B(5, 0.4)$ by computing $P(X = k)$ for $k = 0, 1, 2, 3, 4, 5$.
Verify that the probabilities sum to 1.

\vspace{6pt}
\renewcommand{\arraystretch}{1.8}
\begin{tabularx}{\linewidth}{@{} >{\bfseries}p{1.2cm} X X X X X X @{}}
$k$ & 0 & 1 & 2 & 3 & 4 & 5 \\
\hline
$P(X=k)$ & \blank{1.5cm} & \blank{1.5cm} & \blank{1.5cm} & \blank{1.5cm} & \blank{1.5cm} & \blank{1.5cm} \\
\end{tabularx}

\vspace{6pt}
Sum $= $ \blank{2cm} (should equal 1)

\end{practicebox}

\end{document}
